\section{Respuesta al impulso y al escalón en el modelado}
\label{sec:p1}
\enunciado{¿Por qué es importante la respuesta al impulso y al escalón para modelar sistemas?
Ejemplifique cada una en su proyecto (20\,\%).}

La respuesta al impulso y al escalón son útiles a la hora de modelar sistemas ya que permiten
observar de qué forma reacciona un sistema ante un cambio conocido. Si se mide esta reacción a lo
largo del tiempo, se puede construir un modelo que prediga cómo se comportará el sistema ante
otras condiciones. En el proyecto actual, se aprovechan de la siguiente manera:

\begin{itemize}
  \item \textbf{Respuesta al impulso:} Se aproxima un impulso dando un toque breve al freno
  mientras el disco se encuentra girando y soltarlo inmediatamente después. Este comportamiento
  permite registrar la velocidad antes, durante y después del toque, además de observar el tiempo
  en que tarda en notarse el frenado, así como la velocidad que pierde el disco durante ese breve
  tiempo.
  \item \textbf{Respuesta al escalón:} En este caso se puede generar una respuesta al escalón
  haciendo girar el disco y luego aplicando un frenado y mantener esa presión fija. Esto permite
  como la velocidad responde ante esa presión fija mantenida a lo largo del tiempo.
\end{itemize}

Si se toman los mismos valores iniciales (velocidad) para ambas pruebas, y se comparan con una
prueba donde el disco gira sin frenado, se puede distinguir mejor el efecto que ejerce el frenado
sobre el sistema y por lo tanto generar un modelo del sistema que relacione la relación entre el
frenado y la velocidad del disco.
